\begin{abstract}
Qwen3.5-4B を共通の基盤とする 8 つのファインチューンモデルについて，線形マージ係数を EvoGM 方式の進化的探索で求めた．fitness には，MMLU-Pro，GPQA，BFCL から取った 110 問に対し，チャットテンプレートと思考モードを有効にして\textbf{生成した}回答の素の正答率（上限 4,096 トークン）を用いた．第 1 ラウンド（40 候補，T4 GPU 7.1 時間）で得たマージは，fitness 用とは重複しない 180 問のテストで手動設定の均等マージ Baphomet-1.0-4B（$\lambda_i=0.1375$）と同等の正答率（77.2\% 対 77.8\%）を，半分の生成トークン数で達成した．続いて IRT で高速化した fitness を用いて探索をさらに 2 ラウンド続けた．40 問のサブセットをタスクごとの 2PL モデルで採点し，見込みの薄い候補は途中で打ち切り，有望な候補だけを 110 問で厳密に再評価する方式である（GPU 6.53 時間）．ラウンド 3 の勝者を \textbf{Baphomet-1.1-4B} として公開する．fitness は 0.764 で（ラウンド 1 勝者 0.722，最良の単体ドナー 0.686，Baphomet-1.0 0.603），180 問テストでも数値上は最高だった（79.4\%．Baphomet-1.0 は 77.8\%，基盤モデルは 70.0\%）．ただし Baphomet-1.0 に対する差は有意ではない（不一致ペア 10 対 7，McNemar $p=0.63$）．基盤モデルに対する差は有意である（20 対 3，$p<0.001$）．Baphomet-1.1 の生成トークン数は Baphomet-1.0 より 42\% 少ない（1,834 対 3,180）が，ラウンド 1 勝者（1,498）よりは多い．fitness は温度 0.6 の 1 サンプルであり，勝者はその fitness で選ばれ，テストも小さいため，これらは予備的な結果である．
